% Λύση του 59ου προβλήματος της ενότητας «Άλυτα Προβλήματα».
\subsection*{Λύση Προβλήματος 59 --- ΑΕΠ Ελλάδας και Θεώρημα Μέσης Τιμής}

\begin{alist}
  \item Το μοντέλο \(G\) θεωρείται συνεχές στο \([0,10]\) και
        παραγωγίσιμο στο \((0,10)\). Άρα πληρούνται οι προϋποθέσεις του
        Θεωρήματος Μέσης Τιμής.

        Επομένως υπάρχει τουλάχιστον ένα \(t^*\in(0,10)\) τέτοιο ώστε
        \[
          G'(t^*)
          =
          \frac{G(10)-G(0)}{10-0}.
        \]
        Δηλαδή, κάποια στιγμή μέσα στη δεκαετία, ο στιγμιαίος ρυθμός
        μεταβολής του ΑΕΠ είναι ίσος με τον μέσο ρυθμό μεταβολής της
        δεκαετίας.

  \item Υπολογίζουμε:
        \[
          \frac{G(10)-G(0)}{10}
          =
          \frac{256{,}238-233{,}912}{10}.
        \]
        Άρα
        \[
          \frac{G(10)-G(0)}{10}
          =
          \frac{22{,}326}{10}
          =
          2{,}2326.
        \]
        Επομένως ο μέσος ρυθμός μεταβολής είναι περίπου
        \[
          {
          2{,}233\ \text{δισ. δολάρια ανά έτος}
          }.
        \]

  \item Για το τετραγωνικό μοντέλο
        \[
          G(t)=1{,}51077t^2-12{,}87505t+233{,}91158
        \]
        έχουμε
        \[
          G'(t)=2\cdot1{,}51077t-12{,}87505.
        \]
        Άρα
        \[
          G'(t)=3{,}02154t-12{,}87505.
        \]

        Θέλουμε
        \[
          G'(t^*)=2{,}2326.
        \]
        Επομένως
        \[
          3{,}02154t^*-12{,}87505=2{,}2326.
        \]
        Άρα
        \[
          3{,}02154t^*=15{,}10765
          \quad\Longleftrightarrow\quad
          t^*\approx5.
        \]
        Συνεπώς
        \[
          {t^*=5}.
        \]
        Επειδή \(t=0\) αντιστοιχεί στο \(2014\), το \(t^*=5\)
        αντιστοιχεί στο έτος
        \[
          {2019}.
        \]
\end{alist}
